% Methods snippet: material-realistic PS/PMMA Bragg onions
\subsection{Dispersive PS/PMMA and melanin-like hollow-fill calculations}
The material-realistic calculations used wavelength-dependent complex refractive indices for polystyrene (PS) and poly(methyl methacrylate) (PMMA) measured by Zhang \textit{et al.}\ \cite{Zhang2020PolymerComplexRefractiveIndices}. The tabulated values were linearly interpolated only within their measured range; extrapolation was disabled. Consequently, the material-realistic design set comprised $\lambda_0=470$, 540, 580, 620, 680, 760, and 850~nm, while the idealised 380~nm design was retained only in the constant-index analysis.

For each $\lambda_0$, intact PS/PMMA Bragg onions with PS as material A and the outermost domain were calculated for radial-domain counts up to 120. One intact reference per design wavelength was selected from a colour-performance score combining $\eta_C\eta_L$ and the backward spectral integral. Progressive hollowing was then implemented by replacing consecutive domains from the centre by air while retaining every original radial boundary. The minimum-material geometry was selected subject to $R_{B,\mathrm{integral}}\ge0.90$, visible backward-spectrum correlation $\rho_B\ge0.98$, $\Delta E_{76}\le5$, and $R_{\eta_C\eta_L}\ge0.90$.

For the selected 470 and 580~nm hollow structures, the hollow volume was filled with an idealised melanin-like material having $n_{\mathrm{fill}}=1.75$ and $\kappa_{\mathrm{fill}}\in\{0,0.001,0.003,0.01,0.03,0.1,0.3\}$. The real-index choice is representative of visible-range melanin values reported by Stavenga \textit{et al.}\ \cite{Stavenga2015MelaninRefractiveIndex}; the extinction coefficient was swept because melanin optical constants depend on wavelength, chemistry, packing and preparation \cite{Li2020SyntheticMelaninOpticalConstants}.
